% Generated from traktat/kjelder.json.
\silentchapter{Kjeldenotar}{kjeldenotar}
\begin{widebody}\small
Kjeldetala viser til verk som teksten byggjer på eller vidarefører. Ein kjeldereferanse er ikkje i seg sjølv eit prov for proposisjonen. Tekstlege parallellar er merkte som slike.

\noindent\hypertarget{kjeldenotat-intro-1}{}\textbf{Føreord: Formgrammatikken} Stiny og Gips gjev eit språk for korleis former kan genererast. Samankoplinga med landskap og omhug er vidareføringa i denne teksten.\kjelde{11,18}\par\medskip
\noindent\hypertarget{kjeldenotat-intro-2}{}\textbf{Føreord: teori om designkunnskap} Hatchuel og Weil sin tanke–kunnskapsteori skil mellom konsept og kunnskap. Koplinga til regionane i formrommet er ei vidareføring i FORMLÆRE.\kjelde{16,17,74}\par\medskip
\noindent\hypertarget{kjeldenotat-intro-3}{}\textbf{Føreord: utviklingsbiologi, kybernetikk og agentteori} Desse tradisjonane gjev bakgrunn for målretta system, tilbakekopling og kompetanse på tvers av skalaer. Den samla formmodellen er ikkje eit resultat kjeldene åleine etablerer.\kjelde{4,5,7,21,22}\par\medskip
\noindent\hypertarget{kjeldenotat-intro-4}{}\textbf{Føreord: Omhug} Doctor og medforfattarar drøftar omhug som drivkraft for intelligens; Tronto gjev ein omsorgsetisk bakgrunn. Den avgrensa definisjonen her er teksten si eiga.\kjelde{23,68}\par\medskip
\noindent\hypertarget{kjeldenotat-intro-5}{}\textbf{Føreord: Kapasitet} Arbeida om navigasjonskompetanse er bakgrunn. Skiljet mellom kapasitet, omhug og prediksjon vert formulert og prøvd i FORMLÆRE.\kjelde{21,86}\par\medskip
\noindent\hypertarget{kjeldenotat-1.1}{}\textbf{1.1} tekstleg parallell: Tractatus 1 og 1.1. Opninga omarbeider skiljet mellom verda som det som er tilfelle og som totalitet av ting.\kjelde{108}\par\medskip
\noindent\hypertarget{kjeldenotat-1.203}{}\textbf{1.203} tekstleg parallell: Tractatus 2.021 og 2.024. Substansomgrepet er ein tekstleg parallell; ledd- og nivåargumentet er forma om for denne teksten.\kjelde{108}\par\medskip
\noindent\hypertarget{kjeldenotat-1.2031}{}\textbf{1.2031} tekstleg parallell: Tractatus 2.0211. Regressargumentet omarbeider ein passasje om korleis meining elles ville avhenge av ei anna sann proposisjon.\kjelde{108}\par\medskip
\noindent\hypertarget{kjeldenotat-1.204}{}\textbf{1.204} tekstleg parallell: Tractatus 2.0271. Skiljet mellom det faste objektet og den skiftande konfigurasjonen er omarbeidd til eit skilje mellom analysenivå.\kjelde{108}\par\medskip
\noindent\hypertarget{kjeldenotat-1.21}{}\textbf{1.21} tekstleg parallell: Tractatus 2.032. Definisjonen av struktur har ein nær tekstleg parallell.\kjelde{108}\par\medskip
\noindent\hypertarget{kjeldenotat-1.211}{}\textbf{1.211} tekstleg parallell: Tractatus 2.03. Kjedebiletet er omarbeidd til sambandet mellom delane i ein stol.\kjelde{108}\par\medskip
\noindent\hypertarget{kjeldenotat-1.2123}{}\textbf{1.2123} tekstleg parallell: Tractatus 2.033. Forma som strukturens moglegheit er her vidareført gjennom at strukturen kan gå att i andre objekt.\kjelde{108}\par\medskip
\noindent\hypertarget{kjeldenotat-1.22}{}\textbf{1.22} tekstleg parallell: Tractatus 2.0141. Teksten kallar denne moglegheita valens, der Tractatus nyttar form.\kjelde{108}\par\medskip
\noindent\hypertarget{kjeldenotat-1.221}{}\textbf{1.221} tekstleg parallell: Tractatus 2.0124. Den modal-logiske tanken om at moglege samanhengar er gjevne med objekta, er overført til valensane.\kjelde{108}\par\medskip
\noindent\hypertarget{kjeldenotat-1.223}{}\textbf{1.223} tekstleg parallell: Tractatus 2.0123. Parallellen gjeld sambandet mellom å kjenne objektet og å kjenne samanhengane det kan inngå i.\kjelde{108}\par\medskip
\noindent\hypertarget{kjeldenotat-1.53}{}\textbf{1.53} tekstleg parallell: Tractatus føreordet. Føreordet til Tractatus drøftar problemet med å trekkje ei grense for tanken utan å kunne tenkje begge sider. C/U-formaliseringa er denne teksten si eiga.\kjelde{108}\par\medskip
\noindent\hypertarget{kjeldenotat-7.3}{}\textbf{7.3} tekstleg parallell: Tractatus 6.54. Stigebiletet er ein tydeleg litterær parallell til Tractatus.\kjelde{108}\par\medskip
\noindent\hypertarget{kjeldenotat-7.33}{}\textbf{7.33} tekstleg parallell: Tractatus 6.54. Dette er ei uttrykkeleg vending av stigebiletet: FORMLÆRE held fram lesinga der Tractatus skildrar overskriding.\kjelde{108}\par\medskip
\noindent\hypertarget{kjeldenotat-7.4}{}\textbf{7.4} tekstleg parallell: Tractatus 1. Slutten tek opp att opningsparallellen og fører henne vidare med verknad og navigasjon.\kjelde{108}\par\medskip
\noindent\hypertarget{kjeldenotat-1}{}\textbf{1} historisk bakgrunn: Semper undersøker teknikk, materiale og form. Definisjonen «ordning av skilnader» er formuleringa i denne teksten, ikkje eit dokumentert sitat frå Semper.\kjelde{1}\par\medskip
\noindent\hypertarget{kjeldenotat-1.2}{}\textbf{1.2} omgrepsleg samanlikning: Notata skil klassisk mereologi frå den operasjonelle objektdelinga her. Eit objekt er udeleleg innanfor analysen, ikkje utan vidare fysisk udeleleg.\kjelde{71}\par\medskip
\noindent\hypertarget{kjeldenotat-1.31}{}\textbf{1.31} vidareføring: Raup undersøker formvariasjon i skjell; Mitteroecker og Huttegger drøftar morforom. Den substratuavhengige bruken er ei vidareføring i FORMLÆRE.\kjelde{10,20}\par\medskip
\noindent\hypertarget{kjeldenotat-1.432}{}\textbf{1.432} omgrepsleg bakgrunn: Kauffman-tradisjonen gjev bakgrunnen for det nærliggande moglege. Avgrensinga gjennom regionar og eitt grammatikksteg er formulert i denne teksten.\kjelde{13,59}\par\medskip
\noindent\hypertarget{kjeldenotat-1.5}{}\textbf{1.5} vidareføring: Den materielle delinga er skild frå delinga mellom det kjende, det tenkte og det utenkjelege. Dette er FORMLÆRE si kopling til C-K-teorien.\kjelde{16,17}\par\medskip
\noindent\hypertarget{kjeldenotat-1.6}{}\textbf{1.6} kjeldegrunnlag: Formgrammatikk-tradisjonen ligg til grunn for algebraen og innleiringa. Stiny utviklar dette omgrepsapparatet i Shape: Talking about seeing and doing.\kjelde{11,18}\par\medskip
\noindent\hypertarget{kjeldenotat-1.611}{}\textbf{1.611} kjeldegrunnlag: Å finne delformer gjennom innleiring er sentralt i formgrammatikk. Delformene treng ikkje følgje ei objektdeling som er fastsett på førehand.\kjelde{11,18}\par\medskip
\noindent\hypertarget{kjeldenotat-2.101}{}\textbf{2.101} vidareføring: Michls funksjonalismekritikk er bakgrunn. Påstanden om forma som det som står att etter samtidige trykk, er ei positiv vidareføring i denne teksten.\kjelde{15,77,78}\par\medskip
\noindent\hypertarget{kjeldenotat-2.14}{}\textbf{2.14} vidareføring: Affordanse er relasjonell hos Gibson. Notata koplar henne vidare til landskap og semiotikk; å definere klassen gjennom ei delt konstellasjon av operasjonar er teksten sitt grep.\kjelde{12,55,56,73}\par\medskip
\noindent\hypertarget{kjeldenotat-3.1}{}\textbf{3.1} sameining av tradisjonar: Det adaptive landskapet og det epigenetiske landskapet har ulik historisk bruk. FORMLÆRE fører dei saman i eit formrom; det er ikkje ei uendra overtaking av éin modell.\kjelde{3,8,58}\par\medskip
\noindent\hypertarget{kjeldenotat-3.2}{}\textbf{3.2} vidareføring: Kanalisering er knytt til Waddingtons arbeid om utviklingsrobustheit. Landskapsbiletet vert her brukt på dimensjonar i formrommet.\kjelde{8}\par\medskip
\noindent\hypertarget{kjeldenotat-4.3}{}\textbf{4.3} vidareføring: Arthur er bakgrunn for stiavhengigheit: tidlegare realiseringar endrar vilkåra for seinare val. Bruken på formhistoria er vidareføringa her.\kjelde{14}\par\medskip
\noindent\hypertarget{kjeldenotat-4.4}{}\textbf{4.4} omgrepsleg samanlikning: Nisjekonstruksjon er ein biologisk parallell til at det realiserte endrar vilkåra for dei som kjem etter. Parallellen gjer ikkje alle formdomene biologiske.\kjelde{79}\par\medskip
\noindent\hypertarget{kjeldenotat-5.02}{}\textbf{5.02} vidareføring: Den formelle skildringa viser til Markov-teppe: tilstandar som skil indre og ytre ved vilkårsbestemd uavhengigheit. «Hinne» og sambandet til sjølvet er formulerte i denne teksten.\kjelde{83}\par\medskip
\noindent\hypertarget{kjeldenotat-5.11}{}\textbf{5.11} modellgrunnlag: Fri energi, modell og politikk har grunnlag i desse arbeida. Den presise agentdefinisjonen og korleis han vert brukt på formgjeving, er eit val i FORMLÆRE.\kjelde{19,21,83,84}\par\medskip
\noindent\hypertarget{kjeldenotat-5.113}{}\textbf{5.113} historisk bakgrunn: Kybernetikken gjev den historiske mål–avvik–justering-strukturen. Den probabilistiske omskrivinga og vilkåra for grensetilfellet er teksten sine.\kjelde{4,5,7}\par\medskip
\noindent\hypertarget{kjeldenotat-5.2}{}\textbf{5.2} vidareføring: Doctor og medforfattarar drøftar omhug som drivkraft for intelligens; Tronto gjev ein omsorgsetisk bakgrunn. Skiljet mellom kapasitet, lyskjegle og omhug er formalisert her.\kjelde{23,68}\par\medskip
\noindent\hypertarget{kjeldenotat-5.21}{}\textbf{5.21} operasjonell vidareføring: Kompetanse og navigasjon i rom er kjeldegrunnlaget. Den særskilde terskelmengda \ensuremath{κ}(A) er definert i teksten og må vurderast som del av modellen.\kjelde{21,86}\par\medskip
\noindent\hypertarget{kjeldenotat-5.22}{}\textbf{5.22} vidareføring: Kognitive lyskjegler namngjev grenser for kva mål eit system kan arbeide mot. Den aksevise og tidslege definisjonen her er teksten si presisering.\kjelde{21,22,85,90}\par\medskip
\noindent\hypertarget{kjeldenotat-5.263}{}\textbf{5.263} omgrepsleg bakgrunn: Tilvisinga gjeld prinsipielle grenser for sjølvrepresentasjon. Koplinga til det utenkjelege og til blindskap er vidareføringa i teksten.\kjelde{70}\par\medskip
\noindent\hypertarget{kjeldenotat-5.4}{}\textbf{5.4} omgrepsleg samanlikning: Pyes skilje mellom risiko og vissheit i handverket er ein parallell i notata. Å kalle substratet ein agent av nullte orden er modellgrepet i denne teksten.\kjelde{26,81,92,93}\par\medskip
\noindent\hypertarget{kjeldenotat-5.401}{}\textbf{5.401} modellbakgrunn: Skiljet mellom strukturen og reglane som verkar på han er kopla til eit førehandsmanus om intelligens på fleire skalaer. Kjelda er oppført som preprint.\kjelde{92}\par\medskip
\noindent\hypertarget{kjeldenotat-5.6}{}\textbf{5.6} definisjonsforslag: Liv som oppretthald av lyskjegla er ein definisjon i FORMLÆRE, utvikla med bakgrunn i desse arbeida. Tilvisinga åleine etablerer ikkje definisjonen som allmenn biologisk semje.\kjelde{22,24,85,89}\par\medskip
\noindent\hypertarget{kjeldenotat-5.7}{}\textbf{5.7} syntese i denne teksten: Dei tre kjeldetradisjonane er grunnlag for komponentane. Kjernelikninga er syntesen som vert føreslått i FORMLÆRE, ikkje ei likning henta frå desse kjeldene.\kjelde{11,17,21}\par\medskip
\noindent\hypertarget{kjeldenotat-5.72}{}\textbf{5.72} modellvilkårsvis følgje: Identifikasjonen av formgrammatikk og C-K som grensetilfelle er eit argument i FORMLÆRE. Tilvisingane namngjev tradisjonane som vert samanlikna.\kjelde{11,17}\par\medskip
\noindent\hypertarget{kjeldenotat-6.1}{}\textbf{6.1} vidareføring: Aktørnettverk og distribuert kognisjon er bakgrunn for fleiragentperspektivet. Den bestemte formaliseringa gjennom lyskjegler er teksten si eiga.\kjelde{66,67}\par\medskip
\noindent\hypertarget{kjeldenotat-6.11}{}\textbf{6.11} kjeldegrunnlag: Desse arbeida er bakgrunn for distribuert og fleirskala kompetanse. Overføringa til materiale, handverk og marknad skjer i FORMLÆRE.\kjelde{22,67,85,87}\par\medskip
\noindent\hypertarget{kjeldenotat-6.12}{}\textbf{6.12} biologisk modellbakgrunn: Den biologiske overgangen frå mange agentar til éin er bakgrunn. Vilkåret om delt spenning er teksten si presisering av koplinga.\kjelde{87}\par\medskip
\noindent\hypertarget{kjeldenotat-6.14}{}\textbf{6.14} designfagleg vidareføring: Tilvisingane gjev bakgrunn for distribuert kontroll og samarbeid med kompetente materiale. Koordinering av lyskjegler er teksten si designfaglege formulering.\kjelde{66,67,93}\par\medskip
\noindent\hypertarget{kjeldenotat-6.31}{}\textbf{6.31} omgrepsleg samanlikning: Normans arbeid om daglegdagse ting og bruk er bakgrunn. Stilarten som grensesnitt er ei overføring i FORMLÆRE, ikkje eit direkte sitat.\kjelde{80}\par\medskip
\noindent\hypertarget{kjeldenotat-6.6}{}\textbf{6.6} syntese i denne teksten: Den eksponentielle formmodellen er formulert her. Dei oppførte kjeldene gjev bakgrunn for kompetanse og persistente system, ikkje åleine eit prov for denne fordelinga.\kjelde{86,89}\par\medskip
\noindent\hypertarget{kjeldenotat-7.13}{}\textbf{7.13} omgrepsleg parallell: Notata knyter lesinga på tvers av domene til Batesons mønsterperspektiv. Det er same lesemåte som vert prøvd, ikkje eit prov på ontologisk identitet.\kjelde{72}\par\medskip
\noindent\hypertarget{kjeldenotat-7.15}{}\textbf{7.15} hypotese i denne teksten: Kjeldene gjev omgrep om kompetansegrenser og sjølvrepresentasjon. Den predikerte sviktsignaturen er ein påstand i FORMLÆRE som krev eiga empirisk prøving.\kjelde{21,70}\par\medskip
\end{widebody}
